\documentclass[11pt]{article}

\usepackage[margin=1.15in]{geometry}
\usepackage{amsmath,amssymb,amsthm}
\usepackage{booktabs}
\usepackage{microtype}
\usepackage{xcolor}
\usepackage[colorlinks=true,linkcolor=blue!50!black,citecolor=blue!50!black,urlcolor=blue!50!black]{hyperref}

\newtheorem{theorem}{Theorem}[section]
\newtheorem{lemma}[theorem]{Lemma}
\newtheorem{corollary}[theorem]{Corollary}
\newtheorem{proposition}[theorem]{Proposition}
\newtheorem{remark}[theorem]{Remark}

\newcommand{\N}{\mathbb{N}}
\newcommand{\Q}{\mathbb{Q}}
\newcommand{\R}{\mathbb{R}}
\newcommand{\code}[1]{\texttt{#1}}
\newcommand{\classes}[1]{[3^{#1}]}

\title{An Improved Lower Bound for the $3x+1$ Problem:\\
       $\pi_a(x) \ge x^{0.895}$}
\author{M.~Sharpe\thanks{\texttt{msharpe@irendi.com}. The computations,
  the certification scheme, and this text were produced in collaboration
  with Claude (Anthropic). All final verifications are exact integer
  computations; see Section~\ref{sec:cert}.}}
\date{June 12, 2026; revised August 29, 2026}

\begin{document}
\maketitle

\begin{abstract}
For $a \not\equiv 0 \pmod 3$ let $\pi_a(x)$ count the integers
$n \le x$ whose forward orbit under the $3x+1$ function contains $a$.
Krasikov and Lagarias (2003) proved $\pi_a(x) \ge x^{0.84}$ for all
sufficiently large $x$, by exhibiting a feasible solution to an
explicit linear program $L_k^{NT}(\lambda)$ over the
$3^{k-1}$ congruence classes $m \equiv 2 \pmod 3$ of modulus $3^k$, at
level $k = 11$; this has remained the best published exponent. We
improve the bound to
\[
\pi_a(x) \;\ge\; x^{179/200} \;=\; x^{0.895}
\qquad \text{for all sufficiently large } x,
\]
using their theorem verbatim and new feasible solutions at level
$k = 17$ ($43{,}046{,}721$ classes). The computational unlock is
structural: for fixed $\lambda$ the program is a feasibility question
$c \le F_\lambda(c)$ for a monotone, homogeneous, concave map, so a
feasible point is found by nonlinear power iteration
(Collatz--Wielandt) rather than linear programming, in time linear in
the number of classes per sweep. Soundness does not rest on floating
point: for each claimed exponent $p/q$ we exhibit an integer vector
and verify every inequality of $L_k^{NT}(2^{p/q})$ in exact integer
arithmetic, with the transcendental coefficients replaced by rational
lower bounds certified via integer power comparisons. The method
reproduces the complete published table of Krasikov--Lagarias
($2 \le k \le 11$) to all seven printed decimals, and yields certified
exponents $0.853, 0.863, 0.8724, 0.8812, 0.888, 0.895$ at
$k = 12, \dots, 17$. Separately, the record itself is now MACHINE-VERIFIED: a Lean~4
formalization of the Krasikov--Lagarias inequality system --- including
its ``advanced'' term, which their proof handles by an elimination
argument and which we handle by a strong induction over roots on the
measure $10t + \lfloor \log_2 a^{498}\rfloor$ --- together with a
$177{,}147$-class certificate decided by the Lean kernel, gives
$\pi_1(x) = \Omega(x^{\gamma})$ with
$\gamma = 50\log_2(101182/100000) = 0.8476\ldots > 0.84$ at level
$k = 12$ and $\gamma = 0.8569\ldots$ at $k = 13$, the first
kernel-checked improvements of the 2003 bound
(Section~\ref{sec:formal}).
\end{abstract}

\section{Introduction}

The $3x+1$ function is
$T(n) = n/2$ for even $n$ and $T(n) = (3n+1)/2$ for odd $n$; the
$3x+1$ conjecture asserts that every $n \ge 1$ eventually reaches $1$
under iteration \cite{lagarias-book, lagarias-overview}. For a positive
integer $a \not\equiv 0 \pmod 3$, let
\[
\pi_a(x) \;:=\; \#\{\, 1 \le n \le x \;:\; T^{(j)}(n) = a
\text{ for some } j \ge 0 \,\}.
\]
A sequence of results, beginning with Crandall \cite{crandall} in 1978,
established lower bounds $\pi_1(x) > x^{\beta}$ for increasing
exponents $\beta$: the difference-inequality method introduced by
Krasikov \cite{krasikov} gave $x^{0.43}$, Wirsching \cite{wirsching}
$x^{0.48}$, the nonlinear-programming approach of Applegate and
Lagarias \cite{applegate-lagarias} $x^{0.81}$, and Krasikov--Lagarias
\cite{kl} reached
\[
\pi_a(x) \;\ge\; x^{0.84}
\qquad (a \not\equiv 0 \bmod 3,\; x \ge x_0(a)),
\]
by a computation, performed in 2002 at level $k = 11$, on a linear
program family $L_k^{NT}(\lambda)$ recalled below. This exponent has
remained the best published bound: the 2021 overview
\cite{lagarias-overview} cites it as the state of the art, and we are
not aware of any later improvement. (Tao's celebrated result that
almost all orbits attain almost bounded values \cite{tao} concerns
logarithmic density of a different event and neither implies nor is
implied by pointwise bounds on $\pi_a$.)

Our result is:

\begin{theorem}\label{thm:main}
For every positive integer $a \not\equiv 0 \pmod 3$ there is
$x_0(a)$ such that
\[
\pi_a(x) \;\ge\; x^{179/200} \;=\; x^{0.895}
\qquad \text{for all } x \ge x_0(a).
\]
\end{theorem}

The theorem follows from the Krasikov--Lagarias framework
\emph{verbatim} --- their Theorem~2.2 converts any feasible solution
of $L_k^{NT}(\lambda)$ into the bound with exponent
$\gamma = \log_2 \lambda$, and they prove feasibility lifts from $k$
to $k+1$, so larger $k$ never hurts --- together with a new feasible
solution at $k = 17$, over $43{,}046{,}721$ congruence classes, with
$\lambda = 2^{179/200}$. The contribution of this paper is therefore
computational, in two parts that we believe are of independent
methodological interest:

\begin{enumerate}
\item \textbf{Feasibility without linear programming
(Section~\ref{sec:method}).} For fixed $\lambda$ the program
$L_k^{NT}(\lambda)$ asks for a positive bounded vector $c$ with
$c \le F_\lambda(c)$, where $F_\lambda$ is monotone, positively
homogeneous, and concave (each coordinate is a linear term plus a
minimum of three coordinates, with positive coefficients). Any vector
satisfying $\min_m F_\lambda(c)_m / c_m \ge 1$ is itself a feasible
point --- this is the easy direction of the Collatz--Wielandt
characterization for monotone homogeneous maps (the coincidence of
names is pleasant but accidental: Lothar Collatz of the eigenvalue
bound is Lothar Collatz of the conjecture). Normalized power iteration
on $F_\lambda$ produces such vectors whenever $\lambda$ is below the
feasibility supremum $\lambda_k$. Each sweep costs $O(3^k)$, which is
what moves the frontier from $3^{10}$ classes (LP, 2002 hardware) to
$3^{16}$ and beyond.
\item \textbf{Exact certification (Section~\ref{sec:cert}).} Floating
point is used only to \emph{find} candidate vectors. Each claimed
exponent is a rational $p/q$; the three coefficients of
$L_k^{NT}(2^{p/q})$ are bounded below by dyadic rationals whose
correctness is an integer power comparison; and the candidate vector,
scaled and floored to integers, is checked against every inequality in
exact integer arithmetic --- all $43{,}046{,}721$ of them at $k = 17$.
Since the coefficients enter the right-hand sides positively, rounding
them down only strengthens the verified system; soundness is
one-directional throughout.
\end{enumerate}

As validation, the same machinery reproduces every entry of the
published Krasikov--Lagarias table ($2 \le k \le 11$) to all seven
printed decimal places, including their record point
$\lambda_{11} = 1.7922310$ (Section~\ref{sec:results}).

\section{The Krasikov--Lagarias program}\label{sec:kl}

We recall the objects of \cite{kl}, to which we refer for proofs.
Throughout $\alpha := \log_2 3 = 1.58496\ldots$ and
\[
\classes{k} \;:=\; \{\, m \bmod 3^k \;:\; m \equiv 2 \!\!\pmod 3 \,\},
\qquad |\classes{k}| = 3^{k-1}.
\]
For $m \in \classes{k}$ define
$\phi_k^m(y) := \inf\{ \pi^*_a(2^y a) : a \equiv m \bmod 3^k,\,
a \text{ not in a cycle}\}$, where $\pi^*_a$ is the variant of
$\pi_a$ counting orbits that stay $\le x$ (see \cite[\S2]{kl}). These
functions satisfy a system of difference inequalities
$\mathcal{I}_k$ \cite[Prop.~2.1]{kl}, and Krasikov--Lagarias associate
to it the linear program family
\[
L_k^{NT}(\lambda):\quad \text{minimize } C_k^{\max} \text{ subject to}
\]
\begin{align}
&\text{(L0)}\quad 1 \le c_k^m \le C_k^{\max}
&& m \in \classes{k}, \notag\\
&\text{(L1)}\quad c_k^m \le c_k^{4m} \lambda^{-2}
   + c_{k-1}^{(4m-2)/3} \lambda^{\alpha - 2}
&& m \equiv 2 \!\!\pmod 9, \notag\\
&\text{(L2)}\quad c_k^m \le c_k^{4m} \lambda^{-2}
&& m \equiv 5 \!\!\pmod 9, \notag\\
&\text{(L3)}\quad c_k^m \le c_k^{4m} \lambda^{-2}
   + c_{k-1}^{(2m-1)/3} \lambda^{\alpha - 1}
&& m \equiv 8 \!\!\pmod 9, \notag\\
&\text{(L4)}\quad c_{k-1}^{u} \le c_k^{u},\;\;
   c_{k-1}^{u} \le c_k^{u + 3^{k-1}},\;\;
   c_{k-1}^{u} \le c_k^{u + 2\cdot 3^{k-1}}
&& u \in \classes{k-1}, \notag
\end{align}
with indices taken modulo $3^k$ (note $4m$, $(4m-2)/3$ and $(2m-1)/3$
remain $\equiv 2 \bmod 3$). Their main results, which we use as black
boxes:

\begin{theorem}[{\cite[Thm.~2.2 and \S6]{kl}}]\label{thm:kl}
If $1 \le \lambda \le 2$ and $L_k^{NT}(\lambda)$ has a feasible
solution with principal variables $\{c_k^m\}$, then
$\phi_k^m(y) \ge \Delta_1\, c_k^m\, \lambda^y$ for all $y \ge 0$, with
$\Delta_1 = (4 \max_m c_k^m)^{-1}$; consequently, for every
$a \not\equiv 0 \pmod 3$,
$\pi_a(x) \ge x^{\log_2 \lambda}$ for all sufficiently large $x$.
\end{theorem}

\begin{proposition}[{\cite[\S6]{kl}}]\label{prop:lift}
If $L_k^{NT}(\lambda)$ is feasible, so is $L_{k+1}^{NT}(\lambda)$
(set $c_{k+1}^{m + j\cdot 3^k} := c_k^m$ for $j = 0, 1, 2$). Hence
$\lambda_k := \sup\{\lambda : L_k^{NT}(\lambda) \text{ feasible}\}$ is
non-decreasing in $k$.
\end{proposition}

It is an open problem whether $\lambda_k \to 2$, which would give
$\pi_a(x) \ge x^{1 - \epsilon}$ for every $\epsilon > 0$ \cite[\S6]{kl}.

\section{Feasibility as a nonlinear eigenproblem}\label{sec:method}

Fix $\lambda$ and eliminate the auxiliary variables: given any
principal vector $c = (c_k^m)_{m \in \classes{k}}$, the optimal choice
in (L4) is
$\bar c_{k-1}^{u} := \min\{c_k^u,\, c_k^{u+3^{k-1}},\,
c_k^{u+2\cdot 3^{k-1}}\}$ \cite[eq.~(2.15)]{kl}. Substituting into
(L1)--(L3) defines the map $F_\lambda : \R_{>0}^{\classes{k}} \to
\R_{>0}^{\classes{k}}$,
\[
F_\lambda(c)_m \;=\;
\begin{cases}
\lambda^{-2} c_{4m} + \lambda^{\alpha-2}\, \bar c_{(4m-2)/3}
  & m \equiv 2 \pmod 9,\\[2pt]
\lambda^{-2} c_{4m}
  & m \equiv 5 \pmod 9,\\[2pt]
\lambda^{-2} c_{4m} + \lambda^{\alpha-1}\, \bar c_{(2m-1)/3}
  & m \equiv 8 \pmod 9,
\end{cases}
\]
and $L_k^{NT}(\lambda)$ is feasible if and only if there is a vector
$c$ with $1 \le c_m$, bounded, and $c \le F_\lambda(c)$ pointwise
(then $C^{\max} := \max c$). The map $F_\lambda$ is monotone
($c \le c' \Rightarrow F(c) \le F(c')$), positively homogeneous
($F(tc) = tF(c)$ for $t > 0$), and concave; all coefficients are
positive.

\begin{lemma}\label{lem:cert}
If $c > 0$ satisfies $\min_{m} F_\lambda(c)_m / c_m \ge 1$, then
$c / \min_m c_m$ is a feasible solution of $L_k^{NT}(\lambda)$.
\end{lemma}

\begin{proof}
$c \le F_\lambda(c)$ is the displayed condition; homogeneity preserves
it under rescaling; the rescaled vector has $\min = 1$ and is finite,
so (L0) holds with $C^{\max} = \max c / \min c$; (L1)--(L4) hold with
$\bar c$ as auxiliaries by construction.
\end{proof}

This is the easy (verification) direction of the Collatz--Wielandt
formula for monotone homogeneous maps \cite{gaubert-gunawardena}; the
useful empirical fact, standard in nonlinear Perron--Frobenius theory,
is that normalized power iteration $c \leftarrow F_\lambda(c) /
\|F_\lambda(c)\|_\infty$ converges to vectors achieving
$\min_m F(c)_m / c_m$ close to the spectral value, so that whenever
$\lambda < \lambda_k$ the iteration eventually exhibits a certificate
in the sense of Lemma~\ref{lem:cert}. We stress that no claim in this
paper depends on convergence theory: the iteration is only a search
heuristic, and Lemma~\ref{lem:cert} is checked directly on its output.

Three practical devices matter at scale. (i) Each sweep is two gather
operations and one three-way minimum over arrays of length $3^{k-1}$;
$k = 17$ sweeps take under a second vectorized. (ii) Bisection on
$\lambda$ re-uses the certificate of the last feasible $\lambda$ as a
warm start. (iii) For the largest $k$ we skip bisection entirely: by
Proposition~\ref{prop:lift} the certificate from level $k-1$ lifts to
level $k$, and a single targeted run at a chosen $\lambda$ either
certifies it or not --- any certificate is final, so nothing is lost
by aiming directly at a target exponent.

\section{Exact certification}\label{sec:cert}

Let the claimed exponent be $\gamma = p/q \in \Q$ and
$\lambda = 2^{p/q}$, so $\log_2 \lambda = p/q$ exactly. The three
coefficients of $F_\lambda$ are then $q$-th roots of explicit
rationals:
\[
\lambda^{-2} = \Bigl(\tfrac{1}{2^{2p}}\Bigr)^{1/q}, \qquad
\lambda^{\alpha-2} = \Bigl(\tfrac{3^p}{2^{2p}}\Bigr)^{1/q}, \qquad
\lambda^{\alpha-1} = \Bigl(\tfrac{3^p}{2^{p}}\Bigr)^{1/q}.
\]

\begin{lemma}\label{lem:round}
Let $u/v \le (A/B)^{1/q}$ with $u, v, A, B$ positive integers, which
is equivalent to the integer inequality $u^q B \le v^q A$. If the
system obtained from $L_k^{NT}(2^{p/q})$ by replacing each coefficient
with such a rational lower bound admits a solution, then so does
$L_k^{NT}(2^{p/q})$ itself.
\end{lemma}

\begin{proof}
All coefficients occur positively on right-hand sides of
$\le$-constraints; decreasing them only shrinks the right-hand sides.
\end{proof}

The verifier therefore proceeds as follows. Rational lower bounds with
$v = 2^{64}$ are computed for the three coefficients and certified by
the integer comparison of Lemma~\ref{lem:round} (computed once; the
integers involved have a few hundred thousand bits for our $q \le
2500$). The floating-point vector $c$ found in
Section~\ref{sec:method} is scaled by $2^{48}$, floored to an integer
vector $C$, and clamped below at $2^{48}$ (so $c = C/2^{48} \ge 1$).
Each constraint is then checked as a single integer inequality; for
example (L3) becomes
\[
C_m \cdot v_2 v_C \;\le\; u_2 v_C \cdot C_{4m}
 \;+\; u_C v_2 \cdot \min\bigl(C_{u}, C_{u + 3^{k-1}},
 C_{u + 2\cdot 3^{k-1}}\bigr),
\]
with $(u_2, v_2)$ and $(u_C, v_C)$ the certified bounds for
$\lambda^{-2}$ and $\lambda^{\alpha-1}$. A run that checks all
$3^{k-1}$ constraints yields, by Lemmas \ref{lem:cert} and
\ref{lem:round} and Theorem~\ref{thm:kl}, the bound
$\pi_a(x) \ge x^{p/q}$. No floating-point number survives into the
verified statement.

\section{Results}\label{sec:results}

\textbf{Validation.} Bisection with the iteration of
Section~\ref{sec:method} reproduces the published table of
\cite[Table~2]{kl} for $2 \le k \le 11$ at all seven printed decimals;
in particular the record entry $\lambda_{11} = 1.7922310$ is matched
exactly. \medskip

\noindent\textbf{New levels.} Beyond $k = 11$:

\begin{center}
\begin{tabular}{rrlll}
\toprule
$k$ & classes $3^{k-1}$ & $\lambda$ certified & $\gamma = \log_2\lambda$
& exact certificate $p/q$ \\
\midrule
12 & 177{,}147    & 1.8064235 & 0.8531361 & $853/1000 = 0.853$ \\
13 & 531{,}441    & 1.8188238 & 0.8630058 & $863/1000 = 0.863$ \\
14 & 1{,}594{,}323  & 1.8307723 & 0.8724524 & $2181/2500 = 0.8724$ \\
15 & 4{,}782{,}969  & 1.8419683 & 0.8812483 & $2203/2500 = 0.8812$ \\
16 & 14{,}348{,}907 & $\ge 1.8506089$ & $\ge 0.888$ & $111/125 = 0.888$ \\
17 & 43{,}046{,}721 & $\ge 1.8596099$ & $\ge 0.895$ & $179/200 = 0.895$ \\
\bottomrule
\end{tabular}
\end{center}

\noindent
For $k \le 15$ the third column is the certified end of a bisection
(the true $\lambda_k$ exceeds it by less than $2\times 10^{-7}$); for
$k = 16, 17$ it is a single targeted certificate as described in
Section~\ref{sec:method}, and the true $\lambda_k$ is somewhat larger.
In every row the final column was verified by the exact integer
procedure of Section~\ref{sec:cert}; at $k = 17$ this comprises
$43{,}046{,}721$ verified inequalities.

\begin{proof}[Proof of Theorem~\ref{thm:main}]
The exact certificate at $k = 17$ exhibits a feasible solution of
$L_{17}^{NT}(2^{179/200})$ by Lemmas \ref{lem:cert} and
\ref{lem:round}. Apply Theorem~\ref{thm:kl}.
\end{proof}

\section{Machine verification of the record}\label{sec:formal}

The exponent $0.895$ of Section~\ref{sec:results} rests on
Theorem~2.2 of \cite{kl} and an exact integer check of the certificate
in Python. This section reports a fully machine-verified exponent
above the 2003 record: everything from the dynamics of $T$ to the
final inequality is a Lean~4 proof, with the certificate checked by
the Lean kernel itself.

\subsection*{The obstacle: the advanced term}

Krasikov's difference inequalities for the class infima
$\phi^m_k(y) = \inf\{\pi^*_a(2^y a) : a \equiv m \ (3^k)\}$ read
(\cite{kl}, Prop.~2.1; $\alpha = \log_2 3$)
\begin{align*}
m \equiv 2\ (9):&\quad \phi^m(y) \ge \phi^{4m}(y-2) + \bar\phi^{(4m-2)/3}(y+\alpha-2),\\
m \equiv 5\ (9):&\quad \phi^m(y) \ge \phi^{4m}(y-2),\\
m \equiv 8\ (9):&\quad \phi^m(y) \ge \phi^{4m}(y-2) + \bar\phi^{(2m-1)/3}(y+\alpha-1),
\end{align*}
where $\bar\phi$ is the minimum over the three lifts of a class of
modulus $3^{k-1}$. The last term is \emph{advanced}: the child
$(2a-1)/3$ is smaller than $a$, so its tree is counted at a larger
relative scale $y + \alpha - 1$. A feasible solution $c$ of the linear
program with coefficients $\lambda^{-2}, \lambda^{\alpha-2},
\lambda^{\alpha-1}$ yields $\phi^m(y) \ge \Delta c_m \lambda^y$
(their Theorem~2.2), but a straightforward induction on $y$ cannot
invoke the hypothesis at $y + \alpha - 1 > y$. Krasikov and Lagarias
prove Theorem~2.2 through a back-substitution procedure that
eliminates advanced terms (their \S3--5), with a deletion rule
justified by monotonicity and termination proved via K\H{o}nig's
lemma; the derived system has nested minimizations of depth $> 226$
already at $k = 5$. Truncating the advanced term instead, i.e.\ using
monotonicity to retard it below $y$, is what Applegate--Lagarias did
in 1995 and what the earlier formalization \code{Krasikov50.lean}
does; we computed that the truncated system's exponent saturates near
$0.66$ at every level, while the exact system reaches $0.849$ at
$k = 12$: the advanced term is worth about $0.19$ in the exponent, and
cannot be dispensed with.

\subsection*{A root induction instead of elimination}

We work on the $1/50$ grid of \code{Krasikov50.lean}: caps
$\mathrm{cap}(t) = 2^{\lfloor t/50\rfloor} r[t \bmod 50]$ with a fixed
integer table $r \approx 10^4 \cdot 2^{i/50}$, rate $\mu = p/q$ per
grid step, so that $\lambda = \mu^{50}$ and every certificate
condition is an integer inequality. The three inequalities become,
with offsets in grid steps,
\[
\phi(m, t+100) \;\ge\; \phi(4m, t) \;+\;
\begin{cases}
\bar\phi\big(\tfrac{4m-2}{3}, t + 79\big) & m \equiv 2\ (9),\\[2pt]
0 & m \equiv 5\ (9),\\[2pt]
\bar\phi\big(\tfrac{2m-1}{3}, t + 129\big) & m \equiv 8\ (9),
\end{cases}
\]
the offsets $79 = 100 - 21$ and $129 = 100 + 29$ being the grid
roundings of $2 + (\alpha - 2)$ and $2 + (\alpha - 1)$ (kernel-checked
table conditions $4\,\mathrm{cap}(t) \le 3\,\mathrm{cap}(t+21)$ and
$2\,\mathrm{cap}(t+29) \le 3\,\mathrm{cap}(t)$; the earlier file used
$99$ in place of $129$, which is the truncation).

The key observation is that the advanced term can be handled by a
plain strong induction provided one inducts over \emph{roots} $a$
rather than classes, because root size is a second well-founded
quantity that decreases exactly where the scale increases. Define
\[
M(t, a) \;=\; 10\,t + \lfloor \log_2 a^{498} \rfloor .
\]
Along the three branches, with $b = (2a-1)/3$:
\begin{align*}
(4a,\ t-100):&\quad 10t - 1000 + \lfloor\log_2 (4a)^{498}\rfloor
  = 10t + \lfloor\log_2 a^{498}\rfloor + 996 - 1000,\\
(2b,\ t-21):&\quad \lfloor\log_2 (2b)^{498}\rfloor \le
  \lfloor\log_2 a^{498}\rfloor + 207 \ (\text{as } 3\cdot 2b \le 4a),
  \quad -210 + 207,\\
(b,\ t+29):&\quad \lfloor\log_2 b^{498}\rfloor \le
  \lfloor\log_2 a^{498}\rfloor - 291 \ (\text{as } 3b \le 2a),
  \quad +290 - 291,
\end{align*}
so $M$ drops by $4$, $3$ and $1$ respectively; the two inequalities on
logarithms both reduce to the single numeric fact $3^{498} > 2^{789}$
($498 \log_2 3 = 789.3$). The induction then goes through verbatim:
the certificate inequality at the class of $a$, the three per-root
difference inequalities (which are the internals of the class
inequalities, with the child's class now known exactly, so the minimum
over lifts is replaced by the certificate value at that class), and
the induction hypothesis at the three children. Two features of the
grid are essential. The exact advanced shift $\alpha - 1 =
\log_2(3/2)$ would make the measure degenerate --- the conditions on
the exponent $N$ in $\lfloor \log_2 a^N\rfloor$ become $2N < 100\beta$
and $N \log_2(3/2) > 29.248\,\beta$, i.e.\ $N < 50\beta$ and $N >
50\beta$ --- which is precisely why Krasikov--Lagarias needed
elimination; the grid offset $29 < 50\log_2(3/2) = 29.248$ opens the
window $49.58\beta < N < 50\beta$ (we use $\beta = 10$, $N = 498$), at
a cost of $1/50$ of a doubling on that branch. And inducting over
roots removes the minimum over lifts from the induction and the need
for nonemptiness witnesses for the $3^{k-1}$ classes.

\begin{theorem}[\code{G50.growth\_root}]
Let $k \ge 2$, $1 \le q \le p$, and $c : \N \to \N$ with
$1 \le c_m \le C_{\max}$ on the classes $m \equiv 2\ (3)$ of modulus
$3^k$, satisfying
\begin{align*}
m \equiv 5\ (9):&\quad c_m p^{100} \le c_{4m} q^{100},\\
m \equiv 2\ (9):&\quad c_m p^{100} \le c_{4m} q^{100} + \bar c_{(4m-2)/3}\, p^{79} q^{21},\\
m \equiv 8\ (9):&\quad c_m p^{100} q^{29} \le c_{4m} q^{129} + \bar c_{(2m-1)/3}\, p^{129}.
\end{align*}
Then for every root $a \ge 3$, $a \equiv 2\ (3)$, whose orbit reaches
$8$, and every $t$,
\[
c_{a \bmod 3^k}\; p^t q^{100} \;\le\; \pi^*_a\big(\mathrm{cap}(t)\, a\big)\cdot
\big(C_{\max}\, q^t p^{100}\big).
\]
\end{theorem}

\subsection*{The certificate and the kernel check}

At $k = 12$ ($177{,}147$ classes) the grid system admits
$\mu = 101182/100000$, i.e.\ $\lambda = \mu^{50} = 1.7995$ and
$\gamma = 50 \log_2 \mu = 0.8476$, with an integer certificate of
$32$-bit entries found by power iteration on the monotone map of
Section~\ref{sec:method} (exponents $-100, -21, +29$ in place of
$-2, \alpha-2, \alpha-1$) and checked exactly in Python. In Lean the
certificate is $81$ packed natural-number literals of $2187$ entries
each, read by kernel-accelerated shift and mask; the per-class
condition is a Boolean function built from \code{Nat.ble}, evaluated
over each shard by a divide-and-conquer range checker of logarithmic
depth with a once-proved soundness lemma, and discharged by
\code{decide +kernel} --- the kernel's own reduction, no
\code{native\_decide}, no trusted compiler --- in about fifteen minutes
and $18$~GB. The final theorem:

\begin{theorem}[\code{K12.density\_bound}]
For every $y \ge 0$,
\begin{multline*}
554204 \cdot 101182^{50y} \cdot 100000^{100} \\
\le\; \#\{n \le 80000\cdot 2^y : \text{the orbit of } n \text{ reaches } 1\}
\cdot \big(16777216 \cdot 100000^{50y}\cdot 101182^{100}\big),
\end{multline*}
i.e.\ $\pi_1(x) = \Omega(x^{\gamma})$ with $\gamma = 50\log_2(101182/100000)
= 0.84763\ldots$.
\end{theorem}

Axioms: \code{propext}, \code{Classical.choice}, \code{Quot.sound}
only; no \code{sorry}. This is the first machine-verified exponent
above $0.84$ --- the published record since 2003 --- and, we believe,
the first machine verification of any result in the
Krasikov--Lagarias program at full strength. The grid rounding costs
$0.0055$ against the exact linear program ($0.8531$ at $k = 12$); a
finer grid recovers it. The same files instantiate at any level: the
kernel time grows a little faster than linearly in the number of
classes, so $k = 13$ ($0.8569$) is about an hour and $k = 14$ ($0.870$)
a few hours; the informal $0.895$ at $k = 17$ is the same system at a level
the kernel has not yet been asked to check. The $k = 13$ instance has
also been run: \code{K13.density\_bound} (files \code{KL13Data},
\code{KL13Check0--2}, \code{KL13}; $531{,}441$ classes in $243$ shards,
$2$\,h\,$15$\,min of kernel time) gives $\gamma = 50\log_2(101195/100000)
= 0.8569\ldots$, with the analogous displayed inequality
($387326 \cdot 101195^{50y}\cdot 100000^{100} \le \#\{\cdots\}\cdot
(16777216 \cdot 100000^{50y} \cdot 101195^{100})$).

\subsection*{Trust base, stated precisely}

What is verified: the definitions of $T$, the capped backward tree and
its count, the three difference inequalities for every root, the
growth theorem, the certificate conditions for all $177{,}147$
classes, and the final density inequality. What is not verified and
need not be: how the certificate was found (power iteration in
floating point, then rounded). What is imported: nothing beyond
Lean's kernel and Mathlib. What is claimed: exactly the displayed
inequality, for every $y$, with no ``sufficiently large''.

\section{Discussion}\label{sec:disc}

The marginal exponent gains per level,
$+0.0114, +0.0099, +0.0094, +0.0088$ across $k = 11, \dots, 15$, decay
slowly; whether $\lambda_k \to 2$ (equivalently $\pi_a(x) \ge
x^{1-\epsilon}$ for all $\epsilon$) remains the fundamental open
question about this method \cite[\S6]{kl}, and the data is consistent
with it. Each additional level costs a factor of $3$ in memory and
time; $k = 18$ ($129$M classes) is within reach of a workstation, and
$k \approx 20$ of a GPU implementation. The bottleneck is not the
certification (linear scan) but the iteration count near the
feasibility threshold, which warm starts mitigate.

Three remarks on reliability. First, everything rests on
Theorem~\ref{thm:kl}, which is published and peer-reviewed; we add no
asymptotic analysis of our own, and the new content of this paper is
checkable by a referee in exact arithmetic from the certificate files.
Second, each certificate was checked by \emph{two independently
written verifiers}: the original, and a clean-room re-implementation
using only the language's standard library (no shared code, no
numerical libraries), with the constraint system re-derived from
\cite[\S2]{kl}, a different rational bounding scheme for the
coefficients (denominator $10^{21}$ versus $2^{64}$, binary search
versus unit adjustment), and explicit assertions of the congruence
bookkeeping; both verifiers also detect a single corrupted entry
(negative control). Third, the verification procedure is small (under
a hundred lines) and is itself a candidate for formalization: the
certificate check is a finite integer computation of precisely the
kind proof assistants discharge by reflection, and formalizing
Theorem~\ref{thm:kl} --- an induction over difference inequalities ---
would make the exponent fully machine-verified.

\subsection*{Machine-verified exponents, and the critical constant}

That program is underway in the companion repository \cite{companion},
which contains complete Lean~4 formalizations of integer-grid
weakenings of the difference-inequality method at two strengths, both
proved from the dynamics upward (backward-tree combinatorics,
Krasikov's inequalities, kernel-checked certificates; axioms: Lean's
standard \code{propext}, \code{Classical.choice}, \code{Quot.sound}
only --- no \code{sorry}, no extra axioms).

\begin{enumerate}
\item \textbf{Unit grid} (file \code{Krasikov.lean}): caps $2^y \cdot a$,
an $81$-class certificate at modulus $3^5$ with rational rate
$137/100$, giving the machine-checked bound
$\pi_1(x) = \Omega(x^{0.4543})$ (theorem
\code{density\_lower\_bound}).
\item \textbf{$1/50$ grid} (file \code{Krasikov50.lean}): caps follow
the interleaved-geometric sequence
$\mathrm{cap}(t) = 2^{\lfloor t/50\rfloor}\, r[t \bmod 50]$ for a fixed
$50$-entry integer table $r \approx 10^4 \cdot 2^{i/50}$, so that the
per-step rate is the \emph{rational} $\mu = 1261/1250$ and every
certificate condition is a pure integer inequality. The two odd-branch
terms of the difference inequalities retard by only $21/50$
resp.\ $1/50$ of a doubling (the kernel-checked table conditions; the
margin is $2^{21/50} = 1.3382\ldots > 4/3$). A $2187$-class certificate
at modulus $3^8$, decided by the kernel in chunks, together with
per-class witness roots $a = m + 6561q$ ($q \in \{0,1\}$) whose Terras
orbits reach $8$ in at most $146$ kernel-evaluated steps, yields
(theorem \code{K8.density\_bound}):
\[
\pi_1(x) \;=\; \Omega\!\left(x^{\gamma}\right), \qquad
\gamma \;=\; 50 \log_2\!\frac{1261}{1250} \;=\; 0.632009\ldots
\]
\end{enumerate}

The second exponent crosses a meaningful line: it exceeds
$\log 2/\log 3 = 0.630929\ldots$, the critical odd-step density against
which every quantity of this problem is measured (it is the density a
divergent orbit would have to sustain forever, and the constant
governing the cycle Diophantine bounds). The comparison of the two
numbers is numerical rather than structural --- they measure different
things --- but as a milestone it is sharp: the machine-verified count
of integers reaching $1$ now provably grows faster than
$x^{\log 2/\log 3}$. The exponent $0.632$ exceeds every bound published
before Applegate--Lagarias \cite{applegate-lagarias}; the formal pipeline has since been
scaled past the 2003 record itself (Section~\ref{sec:formal}).

\subsection*{Reproducibility}
The search engine, the one-shot driver, both exact verifiers, the
certificate export format (raw little-endian integers with a hashed
sidecar), and instructions are in the repository accompanying this
paper (\code{analysis/kl\_exponent.py}, \code{kl\_oneshot.py},
\code{kl\_certificate.py}, \code{kl\_export\_certificate.py},
\code{kl\_verify\_independent.py}); all runs are deterministic.
Finding the $k = 17$ certificate takes under an hour on a laptop;
verifying it exactly takes minutes per verifier.

\subsection*{Acknowledgments}
The computations, the certification scheme, and this text were
produced in collaboration with Claude (Anthropic). We thank the
authors of \cite{kl} for a paper written so completely that improving
its constant required no new mathematics --- only twenty-three years
of Moore's law and one change of algorithm.

\begin{thebibliography}{10}

\bibitem{applegate-lagarias} D.~Applegate and J.~C.~Lagarias, Density
bounds for the $3x+1$ problem. II. Krasikov inequalities, \emph{Math.
Comp.} \textbf{64} (1995), 427--438.

\bibitem{crandall} R.~E.~Crandall, On the ``$3x+1$'' problem,
\emph{Math. Comp.} \textbf{32} (1978), 1281--1292.

\bibitem{gaubert-gunawardena} S.~Gaubert and J.~Gunawardena, The
Perron--Frobenius theorem for homogeneous, monotone functions,
\emph{Trans. Amer. Math. Soc.} \textbf{356} (2004), 4931--4950.

\bibitem{krasikov} I.~Krasikov, How many numbers satisfy the $3x+1$
conjecture?, \emph{Internat. J. Math. Math. Sci.} \textbf{12} (1989),
791--796.

\bibitem{kl} I.~Krasikov and J.~C.~Lagarias, Bounds for the $3x+1$
problem using difference inequalities, \emph{Acta Arithmetica}
\textbf{109} (2003), 237--258.

\bibitem{lagarias-book} J.~C.~Lagarias (ed.), \emph{The Ultimate
Challenge: The $3x+1$ Problem}, American Mathematical Society, 2010.

\bibitem{lagarias-overview} J.~C.~Lagarias, The $3x+1$ problem: an
overview, arXiv:2111.02635 (2021).

\bibitem{companion} M.~Sharpe, The Collatz conjecture at the critical
line: machine-verified no-go, density, and cycle theorems, companion
paper and repository, 2026.

\bibitem{tao} T.~Tao, Almost all orbits of the Collatz map attain
almost bounded values, \emph{Forum of Mathematics, Pi} \textbf{10}
(2022), e12.

\bibitem{wirsching} G.~J.~Wirsching, \emph{The Dynamical System
Generated by the $3n+1$ Function}, Lecture Notes in Mathematics 1681,
Springer, 1998.

\end{thebibliography}

\end{document}
